\documentclass[11pt,a4paper]{article}
% Print form of SPEC.md, which remains the canonical reference.

\usepackage[utf8]{inputenc}
\usepackage[T1]{fontenc}
\usepackage{lmodern}
\usepackage{amsmath}
\usepackage{booktabs}
\usepackage{hyperref}
\usepackage{xcolor}
\usepackage{geometry}
\geometry{margin=2.5cm}

\hypersetup{colorlinks=true, linkcolor=blue!60!black, citecolor=blue!60!black, urlcolor=blue!60!black}

\title{3DFC Container Specification\\{\large Version 1: normative reference for \texttt{.3dfc} files}}
\author{}
\date{\today}

\begin{document}
\maketitle

\section{Introduction}

This document specifies the 3DFC container format: a single-file, lossless
archive for ordered sequences of PNG frames. It is intended for implementers
of readers, writers, and validators.

The key words ``MUST'', ``MUST NOT'', ``REQUIRED'', ``SHALL'', ``SHALL NOT'',
``SHOULD'', ``SHOULD NOT'', ``RECOMMENDED'', ``MAY'', and ``OPTIONAL'' are to
be interpreted as described in RFC~2119.

3DFC guarantees pixel identity, not file identity: a compliant reader
reproduces exact pixel values (including transparency). It does not preserve
PNG file bytes, ancillary chunks, or metadata outside the pixel buffer.

\section{Conventions}

All multi-byte integers are unsigned little-endian. Integer widths are written
\texttt{u8}, \texttt{u16}, \texttt{u32}, \texttt{u64}. Byte offsets in tables
are zero-based and half-open unless stated otherwise.

Pixels are stored and reconstructed as straight (non-premultiplied) 8-bit
RGBA (NRGBA): four bytes per pixel in R, G, B, A order, rows top to bottom,
with row stride exactly $\text{width} \times 4$ bytes. There is no row padding
and no filtering.

\section{Versions}

There is a single format version, version~1, with a 40-byte header that always
carries the source total (Section~5).

A reader MUST accept version~1 and MUST reject any other version. A writer
MUST emit version~1.

\section{Overall layout}

A file consists of, in order:

\begin{enumerate}
\item Header (40 bytes).
\item Zero or more packets (Section~\ref{sec:packets}). A valid archive holds
at least one packet when the frame count is non-zero.
\item Names section (Section~\ref{sec:names}).
\item Index section (Section~\ref{sec:index}).
\item Footer, exactly the final 16 bytes (Section~\ref{sec:footer}).
\end{enumerate}

Readers locate the names and index sections through the footer, not by
scanning packets.

\section{Header}

\begin{table}[h]
\centering
\begin{tabular}{@{}llll@{}}
\toprule
Bytes & Type & Field & Requirement \\
\midrule
0--3   & bytes & Magic \texttt{"3DFC"} & MUST match exactly \\
4--5   & u16   & Format version (1 or 2) & Others rejected \\
6--7   & u16   & Tile edge in px, $\ge 8$ & Decoder uses it for GoP geometry \\
8--11  & u32   & Image width in pixels & Identical for all frames \\
12--15 & u32   & Image height in pixels & Identical for all frames \\
16--23 & u64   & Frame count & MUST equal names and index counts, and MUST be at least 1 \\
24--27 & u32   & GOP length upper bound & Informative for readers \\
28     & u8    & Flags; bit 0 = non-opaque alpha present & Informative \\
29--31 & ---   & Reserved & Writers MUST write zero; readers MUST ignore \\
32--39 & u64   & Source total: sum of input PNG sizes & Always present \\
\bottomrule
\end{tabular}
\caption{3DFC v2 header layout.}
\label{tab:header}
\end{table}

Bytes 32--39 are absent in version 1.

\section{Packets}
\label{sec:packets}

Each packet is framed as:

\begin{center}
\begin{tabular}{@{}lll@{}}
\toprule
Bytes & Type & Field \\
\midrule
0     & u8  & Packet type: 0 GoI, 1 GoP, 2 FF, 3 Raw \\
1--4  & u32 & Payload length in bytes, $N$ \\
5--(5+N) & bytes & Payload \\
\bottomrule
\end{tabular}
\end{center}

A reader MUST reject a packet whose declared length exceeds $2^{31}-1$ or
overruns the file. Unknown packet types MUST be rejected. The \texttt{size}
recorded in the index for a packet equals $5 + N$.

\section{GoI payload (type 0)}

The payload is a zstd frame which, once decompressed, is exactly
$\text{width} \times \text{height} \times 4$ bytes of NRGBA pixel data in the
layout of Section~2. A GoI packet fully defines one frame and starts a new
decode chain.

\section{GoP payload (type 1)}

The payload is a zstd frame which, once decompressed, has this layout:

\begin{center}
\begin{tabular}{@{}lll@{}}
\toprule
Bytes & Type & Field \\
\midrule
0--3 & u32 & $C$: number of changed tiles \\
4--(4+4C) & $C \times$ u32 & Changed tile indices, ascending row-major order \\
remainder & bytes & Concatenated residuals, tile by tile, row by row \\
\bottomrule
\end{tabular}
\end{center}

Tiles partition the frame into a $\lceil \text{width}/\text{tile} \rceil$ by
$\lceil \text{height}/\text{tile} \rceil$ grid; tile index $i$ covers column
$i \bmod nx$, row $i \div nx$. Edge tiles are clipped to the frame bounds, so
the last column has width $\text{width}-x$ and the last row has height
$\text{height}-y$ when those are smaller than the tile edge.

Residuals are stored per byte as $(\text{current}-\text{previous}) \bmod 256$,
where \texttt{previous} is the fully reconstructed preceding frame in the same
chain. Reconstruction adds each residual byte to the previous frame byte
modulo 256. A compliant decoder MUST verify that the residual section holds
exactly the bytes implied by the listed tile rectangles, and MUST reject short
data.

\section{FF payload (type 2)}

The payload is stored uncompressed (the inner streams are already compressed)
with this layout:

\begin{center}
\begin{tabular}{@{}lll@{}}
\toprule
Bytes & Type & Field \\
\midrule
0--3 & u32 & $F$: frame count of the segment \\
4--7 & u32 & $R$: length of the RGB stream in bytes \\
8--11 & u32 & $A$: length of the alpha stream in bytes (0 if absent) \\
12--(12+R) & bytes & RGB stream: Matroska, lossless H.264 RGB (\texttt{x264rgb}, \texttt{qp~0}) \\
12+R--(12+R+A) & bytes & Alpha stream if $A > 0$: Matroska, lossless H.264 grayscale (\texttt{qp~0}) \\
\bottomrule
\end{tabular}
\end{center}

Decoding yields $F$ RGB frames and, when present, $F$ single-channel frames;
frame $k$ of the segment is recombined as RGB planes plus alpha plane (or
alpha 255 throughout when the alpha stream is absent). Pixel values MUST
round-trip bit-exactly; in particular, no YUV conversion is permitted in
compliant FF segments.

\section{Raw payload (type 3)}

The payload is a complete source PNG file, stored verbatim. On restore it is
written out byte-for-byte, including all ancillary chunks. This packet type
carries 16-bit content and any group the encoder declined to compress.

\section{Names section}
\label{sec:names}

\texttt{count:u64} equal to the header frame count, followed by one entry per
frame in frame order: \texttt{len:u16} (byte length, not NUL-terminated) and
\texttt{len} bytes of UTF-8 filename. Names MUST NOT exceed 65535 bytes.
Names SHOULD be basenames; readers MUST treat them as
untrusted (see Section~\ref{sec:security}).

\section{Index section}
\label{sec:index}

\texttt{count:u64} equal to the header frame count, followed by one 24-byte
record per frame in frame order:

\begin{center}
\begin{tabular}{@{}lll@{}}
\toprule
Bytes & Type & Field \\
\midrule
0--7 & u64 & Byte offset of the packet start (the type byte) \\
8--11 & u32 & Packet size ($5 + $ payload length) \\
12 & u8 & Flags (see below) \\
13--14 & u16 & Position of the frame inside an FF segment & MUST be 0 for non-FF frames \\
15--22 & u64 & Frame hash (Section~\ref{sec:hashes}) \\
23 & --- & Reserved; writers MUST write zero, readers MUST ignore \\
\bottomrule
\end{tabular}
\end{center}

Flag bits:

\begin{description}
\item[\texttt{0x01} Key:] the frame starts an independent decode chain (GoI
packets and the first frame of every FF segment).
\item[\texttt{0x02} FF:] the frame belongs to an FF segment.
\item[\texttt{0x04} Raw:] the frame is a Raw packet.
\end{description}

The first frame of an FF segment carries \texttt{0x01} \texttt{|} \texttt{0x02}.
Every Go run opens with a \texttt{0x01} frame whose packet type is GoI. Raw
records MUST carry \texttt{0x04} alone, without the Key bit.

\section{Hashes}
\label{sec:hashes}

Frame hashes are the first 64 bits (little-endian) of SHA-256 over:

\begin{itemize}
\item the reconstructed NRGBA pixel buffer, for GoI, GoP, and FF frames;
\item the verbatim PNG file bytes, for Raw frames.
\end{itemize}

Hashes are integrity checks, not content identifiers.

\section{Footer}
\label{sec:footer}

The final 16 bytes are \texttt{names\_offset:u64} followed by
\texttt{index\_offset:u64}, both absolute file offsets. A reader MUST verify
that both offsets lie within the file and that
$\text{names\_offset} \le \text{index\_offset}$.

\section{Decode algorithm}

To restore frame $k$:

\begin{enumerate}
\item Read its index record. If it is a Raw frame, copy the payload to the
output path named by the names section and stop.
\item If it is an FF frame, decode the whole segment packet, select position
\texttt{segPos}, recombine, and stop.
\item Otherwise, walk backwards to the nearest record with the Key flag,
decode forward applying GoI/GoP reconstruction, and emit frame $k$.
\end{enumerate}

Decoders MUST validate packet types against index flags and reject
inconsistent archives.

\section{Resume sidecars}

Encoders MAY support interruption and resume with three sidecar files next to
the target path \texttt{out}:

\begin{itemize}
\item \texttt{out.partial}: packet data written so far, beginning with the
header.
\item \texttt{out.resume.idx}: concatenated 24-byte index records for
committed frames.
\item \texttt{out.resume.json}: checkpoint object with \texttt{version},
\texttt{fingerprint}, \texttt{config}, \texttt{next\_segment},
\texttt{data\_offset}, and \texttt{entry\_count}.
\end{itemize}

The fingerprint is hex-encoded SHA-256 over, per source frame in order: name
length (\texttt{u64le}), name bytes, and file SHA-256. The config string binds
the format version, GOP/tile/backend/preset/fps options, and FFmpeg
availability. A resume run MUST verify fingerprint, configuration, header
bytes, and that the partial file and journal are at least as long as the
checkpoint claims, truncating both to the checkpointed lengths before
continuing. Sidecars are removed after a successful atomic commit of the final
file.

\section{Security considerations}
\label{sec:security}

Readers MUST NOT write outside the destination directory: reject absolute
names, path separators, and empty names, and detect case-insensitive
collisions between restored filenames. Readers SHOULD cap decompressed sizes
(e.g.\ against $\text{width} \times \text{height} \times 4$ plus small
overhead) and packet lengths before allocating. The 64-bit truncated hashes
detect corruption, not malicious tampering.

\section{Versioning and compatibility}

Minor, backward-compatible additions use new reserved bytes or new packet
types; readers MUST reject what they do not understand rather than guess. The
version field is bumped only for breaking changes, at which point this
specification is revised and reissued.

\end{document}
